\section{Une septième arche, tirée de la théorie de l'homotopie : les catégories test}\label{sec:homotopical}

Les \emph{catégories test} de Grothendieck sont les petites catégories $A$ dont les préfaisceaux modélisent les types d'homotopie ; $\Delta$ est le premier exemple. Leur caractérisation lit un objet de $\PSh(A)$, le classifiant des sous-objets, de façon homotopique ; on peut donc la faire tourner comme une arche dont l'autre côté est la lecture logique du même objet que donne le pont~\ref{br:dm}. La note associée \cite[section \q{Catégories test et arche De Morgan}]{bridges-notes} le fait sur les questions des dossiers 108, 118 et 134-2 de Grothendieck ; nous en donnons ici la sortie.

Fixons un localisateur fondamental $W$ sur les petites catégories, autre que les deux localisateurs grossiers (par exemple les foncteurs dont le nerf est une équivalence faible). Pour un préfaisceau $X$, notons $A/X$ sa catégorie des éléments. L'\emph{objet de Lawvere} $L_A$ est le préfaisceau des cribles, c'est-à-dire le classifiant des sous-objets de $\PSh(A)$. Une \emph{catégorie test faible} est une catégorie pour laquelle $A/L_A$ est asphérique, entre autres conditions, et une \emph{catégorie test locale} une catégorie dont les tranches $A/a$ sont des catégories test faibles \cite[déf.~4.1.8, 4.1.12]{cisinski-308}.

\begin{bridge}[Catégories test]\label{br:test}
$A$ est une catégorie test locale si et seulement si $L_A$ est localement asphérique : $A/(a\times L_A)$ est $W$-asphérique pour tout objet $a$ \cite[th.~4.1.19]{cisinski-308}, \cite[th.~1.5.6]{maltsiniotis-301}. Lean : non formalisé ; Mathlib n'a pas de localisateurs fondamentaux. Le fichier \texttt{lean/bridges/Bridges/FundamentalLocalizers.lean} de l'édition démontre les conséquences ci-dessous dans un cadre axiomatique.
\end{bridge}

\begin{theorem}\label{th:testore}
Une petite catégorie non vide qui satisfait la condition d'Ore n'est ni une catégorie test faible ni une catégorie test locale. De façon équivalente, le topos des préfaisceaux sur une catégorie test faible ou locale n'est jamais de De Morgan ; pour une catégorie test locale $A$, aucune tranche $\PSh(A)/X$ avec $X$ non vide n'est de De Morgan non plus.
\end{theorem}

\begin{trace}
\item[tablier] $\PSh(A)$ pour une petite catégorie $A$, avec son classifiant des sous-objets $L_A$.
\item[arche] le pont~\ref{br:dm} d'un côté, le pont~\ref{br:test} de l'autre.
\item[spécialiser] la condition d'Ore dit qu'un crible non vide ne se tire jamais en arrière sur le crible vide ; de façon équivalente, les cribles vides se détachent : $L_A = L_A^{\emptyset}\sqcup L_A^{+}$. Dans le langage des topos, cela dit que $\bot\colon 1\to\Omega$ a un complément, ce qui caractérise les topos de De Morgan \cite[\S1]{caramello-demorgan}, \cite[prop.~14]{caramello-chamoun}.
\item[réécrire] alors $A/(a\times L_A)$ est réunion disjointe de deux morceaux non vides (le crible vide et le crible maximal sur $a$), donc n'est pas connexe ; pour $W$ non grossier, une catégorie asphérique est connexe \cite[lemme~1.7.22]{maltsiniotis-301}.
\item[vérifier] l'argument direct est l'étape de réécriture ; Lean : axiomatique seulement (\lean{not_isDeMorganTopos_slice_of_isLocalTest}, dans \texttt{lean/bridges} de l'édition).
\item[nouveauté] non trouvé tel quel ; voir ci-dessous.
\end{trace}

\noindent Les tranches suivent, puisque $A/X$ est une catégorie test locale dès que $A$ l'est \cite[prop.~4.1.14 (v)]{cisinski-308}, \cite[rem.~1.5.4 b)]{maltsiniotis-301}, et que $\PSh(A)/X\simeq\PSh(A/X)$. Ainsi les ensembles simpliciaux, cubiques, les $\Theta$-ensembles et les ensembles dendroïdaux ne forment pas des topos de De Morgan, ni aucune de leurs tranches ; dans $\Delta$, le cospan $[0]\rightrightarrows[1]$ des deux faces ne se complète pas. La réciproque est fausse : $\Delta$ ne satisfait pas la condition d'Ore et est une catégorie test.

\status{} Non trouvé tel quel, ni dans \cite{cisinski-308}, \cite{maltsiniotis-301}, \emph{Pursuing Stacks}, ni par TheoremSearch et alphaXiv : aucun des textes sur les catégories test ne mentionne la condition d'Ore ni la loi de De Morgan, même si \emph{Pursuing Stacks} \S88 dit sans démonstration qu'une catégorie test n'a pas d'objet initial. L'énoncé est une combinaison immédiate de Caramello--Chamoun (topos de De Morgan si et seulement si $\bot$ a un complément), du théorème~4.1.19 de Cisinski et du lemme~1.7.22 de Maltsiniotis. Le résultat publié le plus proche est celui de Mielke : un topos est de De Morgan si et seulement s'il est \emph{localement homotopiquement trivial}, tout intervalle d'une tranche se scindant en deux morceaux qui séparent ses extrémités, la condition d'Ore en étant la forme préfaisceaux \cite[th.~1.5, cor.~1.6]{mielke}. Ses intervalles sont totalement ordonnés et son homotopie n'est pas l'asphéricité de Grothendieck, de sorte que le théorème~\ref{th:testore} n'est pas un cas du sien. Le mécanisme est pourtant le même : dans un topos de De Morgan, l'objet des valeurs de vérité se scinde en deux morceaux.
